\section{Разбиение на документы и выделение текста}

За документ принята одна веб-страница: у Википедии --- статья о фильме, сериале или
человеке, у film.ru --- рецензия, подборка или интервью. Такое разбиение напрашивается
само: каждая сущность на обеих площадках уже живёт по собственному адресу.

Карточки фильмов film.ru документами не считаются. Текста в них --- одна аннотация на
несколько десятков слов, а всё остальное (актёры, рейтинги, расписание показов) ---
структурированные поля, которые лучше брать из метаданных, а не из текста.

Для анализа вручную сохранено 13 страниц: 7 из Википедии и 6 с film.ru. Брались
документы разных типов и разного объёма: в Википедии --- статьи о двух фильмах, две
заглушки, сериал и две биографии; на film.ru --- три рецензии разной длины, подборка,
интервью и одна карточка фильма для сравнения.

Текст выделялся так: берётся содержимое основного контейнера, из него выбрасываются
\texttt{<script>}, \texttt{<style>} и таблицы, снимаются остальные теги, раскрываются
HTML-сущности, схлопываются пробелы.

Способ грубый. Выбрасывая таблицы, мы теряем вместе с версткой и карточку фильма с
полезными полями (год, страна, режиссёр); зато разделы \emph{Примечания} и
\emph{Ссылки} в текст попадают, хотя это библиографический шум. Для оценки корпуса
этого достаточно, при индексации правила придётся уточнять.

\pagebreak
